\section{The Research Context Graph}
\label{sec:research-graph}

A multi-agent scientific system must maintain shared state across its agents: the
hypotheses under consideration, the evidence bearing on each, and the steps that
remain. Existing systems typically hold this state as natural-language
dialogue, relaying intermediate results between agents, or to a central
orchestrator, as
prose~\citep{boiko2023coscientist,google2025coscientist,swanson2025virtuallab}.
Prose is an inadequate coordination substrate for this purpose along three axes.
It is \emph{unstructured}, and therefore cannot be queried---the orchestrator
cannot ask whether any hypothesis is ready for verification without re-reading and
re-interpreting the transcript. It is \emph{unbounded}, growing monotonically with
the length of the investigation, so per-step context cost scales with history
rather than with the task. And it is \emph{unverifiable}: an agent can assert an
arbitrary claim in free text, and consuming agents have no mechanism to
distinguish a licensed result from a fabricated one.

We instead represent shared state as a typed, schema-validated property graph
that constitutes the working memory of a single investigation, rooted at one
research question. The schema instantiates a six-layer meta-model of the
scientific process, comprising research questions, hypotheses, the verification
methods and confirmation criteria that would test them, the evidence and
conclusions they yield, and the tools, resources, and constraints they depend on.
Agents do not exchange results as messages; each writes typed nodes to the shared
graph. The orchestrator reads the graph state to select the next action, and each
worker agent is supplied a \emph{context slice}---the assigned node together with
its one- or two-hop neighbourhood---rather than the full history.

\paragraph{Data model.}
A node is a tuple $\langle \mathit{id}, \mathit{type}, \mathit{attrs},
\mathit{status}, \mathit{source}, \mathit{history} \rangle$; an edge is a typed,
directed relation over an ordered pair of nodes. The type system is fixed and
finite. Each node type defines a status lifecycle: a \textsc{Hypothesis}
transitions
$\mathit{formulated} \rightarrow \mathit{under\_verification} \rightarrow
\{\mathit{confirmed}, \mathit{refuted}, \mathit{postponed}\}$; \textsc{Evidence}
carries a mandatory subtype (literature, experimental, computational, expert, or
meta) and links to the hypotheses it supports, refutes, or refines;
\textsc{Conclusion} nodes are grounded in evidence via \emph{based\_on} edges;
\textsc{VerificationMethod} and \textsc{ConfirmationCriteria} encode how a
hypothesis is to be tested and the threshold for confirmation; and \textsc{Tool},
\textsc{Resource}, and \textsc{Constraint} nodes represent the means and limits of
the study. Edges are range-restricted to admissible type pairs---for example,
\emph{motivates} $(Q \rightarrow H)$, \emph{tested\_by} $(H \rightarrow VM)$, and
\emph{supports} $(E \rightarrow H)$---so the graph is a constrained model of the
reasoning process rather than an unconstrained association network.

\paragraph{Validation at the write boundary.}
Each write is validated against the schema before it is committed: node type and
status must be admissible, a status change must follow the corresponding
lifecycle, and an edge must connect a permitted type pair. Writes are additionally
constrained by a per-role permission table, and the identity of the writing agent
is obtained from the trusted tool-call context rather than passed as an argument,
so it cannot be forged. Consequently, the agent that generates hypotheses may
create \textsc{Hypothesis} nodes but cannot transition them to
\emph{confirmed}; only the verification role holds that permission. The write
boundary thus enforces, in code rather than in prompt text, that every recorded
state transition is one the process licenses---a structural constraint on the
class of states an agent can introduce.

\paragraph{Properties.}
The representation provides four properties that a dialogue transcript does not.
First, \emph{queryable control flow}: because state is typed, the orchestrator
selects actions by evaluating named predicates over the graph---hypotheses whose
required tools are all available (\emph{ready} for verification), evidence that
refutes a still-open hypothesis, or conclusions whose criteria are satisfied
(\emph{closable})---rather than by re-interpreting prose. Second,
\emph{provenance}: each conclusion is traceable through \emph{based\_on} edges to
its supporting evidence and onward to the methods and sources that produced it,
and each node retains a status history recording every transition and its author;
the completed investigation is an inspectable, reproducible artefact. Third,
\emph{bounded context}: supplying each agent a local slice rather than the full
transcript keeps per-step context size a function of the immediate task rather
than of accumulated history, reducing both cost and the error modes associated
with long, undifferentiated context. Fourth, \emph{explicit negative results}:
refuted and postponed hypotheses are retained rather than discarded, falsification
is a defined state transition, and a hypothesis under verification is locked
against concurrent re-entry. Because agents coordinate through validated shared
memory rather than point-to-point dialogue, they are decoupled, which admits
scaling to additional agents and to distributed (agent-to-agent) execution.

The research context graph is the reasoning-level counterpart of the system's
execution trace: the execution trace records the actions the agents took, whereas
the research context graph records the claims the investigation established, under
a schema that keeps that record structured, attributable, and resistant to
fabrication.
